\section{Conclusion}\label{sec:conclusion}

We ask which parts of chart construction carry the return ranking and where that ranking retains economic value. We compare scaling and model choices with recent histories and portfolio rules fixed, then examine stock-level overlap, fusion, tradability, and post-2019 stability.

Local high--low scaling makes a strong weekly ranking accessible to the logistic classifier, while CNN1D on the continuous, locally scaled sequence closely matches the chart CNN. Because temporal convolution contributes mainly at the weekly horizon, the results point to short-lived information in the ordered, normalized price path. The binary image is not necessary to recover the portfolio performance, although the architecture differences prevent a pure test of pixel layout.

The chart CNN, CNN1D, and multimodal scores exhibit substantial but incomplete rank overlap. Similar Sharpe ratios can coexist with imperfect full-sample correlations because the portfolios depend on agreement in the extreme deciles. Neither fusion design consistently improves on the strongest individual model. The residual score differences therefore offer little portfolio-level complementarity, but they do not imply that the representations contain identical information.

The unrestricted equal-weighted portfolios remain strong under the common cost schedule, but net performance turns negative under value weighting and modest tradability screens. Conventional factor adjustment does not absorb the unrestricted spread, yet its concentration in difficult-to-trade stocks makes market impact, borrowing constraints, capacity, and microstructure explanations difficult to dismiss. Because the screens do not reproduce a matched largest-stock exercise, they support a more restrictive assessment under alternative tradability criteria rather than a direct reversal of earlier evidence. The learned models and short-horizon price rules also weaken after 2019, which is consistent with a common price-history component losing strength but does not identify whether crowding, market change, or sampling variation is responsible.

Taken together, the evidence is most consistent with a short-lived, locally normalized price-history ranking rather than a durable image-specific effect. The design establishes that the chart CNN's weekly portfolio performance is recoverable without the binary image and that its measured economic value is concentrated in illiquid stocks; it does not identify the precise behavioral or microstructure mechanism behind the ranking.
